Los resultados confirman que no existen las balas de plata, es decir que no existe un algoritmo que sea infalible, sino una relación de compromiso entre optimalidad y costo computacional. La fuerza bruta garantiza la solución óptima pero su crecimiento exponencial la vuelve infactible para $n > 25$, mientras que la \textit{programación dinámica} logra la misma optimalidad con un costo pseudopolinomial, a cambio de un uso de memoria considerablemente mayor debido al uso de la tabulación. Entre las heurísticas, \textit{greedy ratio} demostró ser la opción más equilibrada, alcanzando de forma consistente --- entre la mayoria de las instancias--- soluciones de calidad óptima con tiempos de ejecución en órdenes de magnitud menores que la \textit{programación dinámica}, mientras que \textit{greedy max}, pese a ser el algoritmo más rápido y liviano en memoria, entrega soluciones de baja calidad, sin una dependencia clara respecto a ninguna de las variables de entrada. En definitiva, la elección del algoritmo debe ajustarse al contexto: la \textit{programación dinámica} cuando la optimalidad es indispensable, y \textit{greedy ratio} cuando se prioriza la eficiencia con una pérdida de calidad aceptable.